% =====================================================================
%  ScholarPath — دليل الشرح العربي/الإنجليزي — الـ preamble المشترك
%  Compile with:  xelatex  (NOT pdflatex — Arabic needs XeTeX)
%  عربي أساسي (RTL) + مصطلحات إنجليزية داخلية + صناديق mdframed
% =====================================================================

% ── اللغة والخطوط ─────────────────────────────────────────────────────
\usepackage{fontspec}
\usepackage{polyglossia}
% numerals=maghrib → أرقام لاتينية (0-9) للأقسام والصفحات، أنسب لمستند تقني ثنائي اللغة
\setmainlanguage[numerals=maghrib]{arabic}
\setotherlanguage{english}
\newfontfamily\arabicfont[Script=Arabic,Scale=1.06]{Amiri}
\newfontfamily\englishfont{TeX Gyre Termes}
\setmainfont{TeX Gyre Termes}
% خط أحادي المسافة للأكواد/المعرّفات
\newfontfamily\monofont{Consolas}[Scale=0.92]

% ── الهندسة والألوان ─────────────────────────────────────────────────
\usepackage[a4paper,margin=2.3cm,top=2.7cm,bottom=2.5cm,headheight=26pt,headsep=12pt]{geometry}
\usepackage[table]{xcolor}
\definecolor{spInk}{HTML}{1F2A44}
\definecolor{spPrimary}{HTML}{2A4D9B}
\definecolor{spAccent}{HTML}{0E7C7B}
\definecolor{spHow}{HTML}{1565C0}
\definecolor{spWhy}{HTML}{2E7D32}
\definecolor{spSec}{HTML}{B23A00}
\definecolor{spQA}{HTML}{5E35B1}
\definecolor{spGrayBg}{HTML}{F4F6FB}
\definecolor{spRule}{HTML}{C9D2E8}
\definecolor{spLightBlue}{HTML}{E3F2FD}
\definecolor{spLightGreen}{HTML}{E8F5E9}
\definecolor{spLightPurple}{HTML}{EDE7F6}
\definecolor{spLightOrange}{HTML}{FBE9E7}
\definecolor{spLightTeal}{HTML}{E0F2F1}
\definecolor{spLightYellow}{HTML}{FFF9C4}
\definecolor{spYellow}{HTML}{F9A825}
\definecolor{spOrange}{HTML}{E65100}

% ── حزم عامة ─────────────────────────────────────────────────────────
\usepackage{booktabs}
\usepackage{array}
\usepackage{tabularx}
\usepackage{longtable}
\usepackage{enumitem}
\usepackage{amsmath}
\usepackage{amssymb}
\usepackage{graphicx}
\usepackage{pdflscape}   % صفحات landscape لإدماج الديجرامات الرسمية العريضة
\usepackage{listings}
\usepackage{tikz}
\usetikzlibrary{arrows.meta,positioning,backgrounds,calc,fit}
% ── أنماط عُقَد الرسومات (تُستخدَم داخل \begin{LTR}) — ألوان متّسقة مع الكتاب ──
\tikzset{
  pbox/.style={rectangle,rounded corners=4pt,text width=2.7cm,minimum height=0.8cm,
    align=center,draw=spPrimary!65,fill=spLightBlue,font=\small\bfseries,text=spPrimary!80},
  sbox/.style={rectangle,rounded corners=3pt,text width=2.7cm,minimum height=0.72cm,
    align=center,draw=spPrimary!40,fill=spGrayBg,font=\small,text=spInk},
  gbox/.style={rectangle,rounded corners=4pt,text width=2.7cm,minimum height=0.8cm,
    align=center,draw=spQA!70,fill=spLightPurple,font=\small\bfseries,text=spQA!80},
  fbox/.style={rectangle,rounded corners=3pt,text width=2.7cm,minimum height=0.72cm,
    align=center,draw=spSec!65,fill=spLightOrange,font=\small,text=spSec!80},
  tbox/.style={rectangle,rounded corners=3pt,text width=2.7cm,minimum height=0.72cm,
    align=center,draw=spAccent!65,fill=spLightTeal,font=\small,text=spAccent!80},
  srcbox/.style={rectangle,rounded corners=4pt,text width=2.6cm,minimum height=0.66cm,
    align=center,draw=spWhy!60,fill=spLightGreen,font=\scriptsize,text=spWhy!80},
  dbbox/.style={rectangle,rounded corners=3pt,text width=2.1cm,minimum height=0.8cm,
    align=center,draw=spAccent!70,fill=spLightTeal,font=\small\bfseries,text=spAccent!80,
    double,double distance=1.4pt},
  dbox/.style={rectangle,rounded corners=2pt,text width=2.6cm,minimum height=0.7cm,
    align=center,draw=spYellow!85!black,fill=spLightYellow,font=\scriptsize\itshape,text=spOrange!90},
  wbox/.style={rectangle,rounded corners=3pt,text width=11.4cm,minimum height=0.68cm,
    align=center,draw=spPrimary!35,fill=spGrayBg,font=\small,text=spInk},
  actor/.style={rectangle,rounded corners=2pt,text width=2.2cm,minimum height=0.7cm,
    align=center,draw=spInk!45,fill=white,font=\small\bfseries,text=spInk},
  arr/.style={->,>=Stealth,semithick,color=spInk!55},
  arrd/.style={->,>=Stealth,semithick,dashed,color=spInk!40},
  arrb/.style={<->,>=Stealth,semithick,color=spInk!55},
  lbl/.style={font=\scriptsize,text=spInk!60},
}
\usepackage[framemethod=default]{mdframed}
\usepackage{titlesec}
\usepackage{fancyhdr}
\usepackage[hidelinks]{hyperref}

% ── أدوات نصية مساعدة ────────────────────────────────────────────────
% مصطلح إنجليزي داخل جملة عربية: اتجاه LTR + خط لاتيني نظيف
\newcommand{\en}[1]{\textenglish{#1}}
% معرّف كود/تقني بخط أحادي المسافة
\newcommand{\code}[1]{{\monofont #1}}

% ── تنسيق العناوين (RTL: الرقم على اليمين) ───────────────────────────
% أرقام الأقسام المركّبة (2.1) تُجبَر على الاتجاه LTR حتى تُقرأ صحيحة لا معكوسة
\renewcommand{\thesection}{\textenglish{\arabic{section}}}
\renewcommand{\thesubsection}{\textenglish{\arabic{section}.\arabic{subsection}}}
\renewcommand{\thesubsubsection}{\textenglish{\arabic{section}.\arabic{subsection}.\arabic{subsubsection}}}
\titleformat{\section}{\Large\bfseries\color{spPrimary}}{\thesection.}{0.5em}{}
\titleformat{\subsection}{\large\bfseries\color{spAccent}}{\thesubsection}{0.5em}{}
\titleformat{\subsubsection}{\normalsize\bfseries\color{spInk}}{\thesubsubsection}{0.5em}{}
\titlespacing*{\section}{0pt}{16pt}{6pt}
\titlespacing*{\subsection}{0pt}{11pt}{4pt}
\setcounter{secnumdepth}{3}
\setcounter{tocdepth}{2}
\setlist{topsep=2pt,itemsep=1pt,parsep=0pt,leftmargin=1.6em,rightmargin=0.2em}
\setlength{\parindent}{0pt}
\setlength{\parskip}{5pt plus 1pt minus 1pt}

% ── الترويسة والتذييل ────────────────────────────────────────────────
\pagestyle{fancy}\fancyhf{}
\renewcommand{\headrulewidth}{0.4pt}
\fancyhead[R]{\small\color{spPrimary}\textbf{ScholarPath}}
% الترويسة اليسرى: «الفصل N» — نتجنّب عرض عنوان القسم في الترويسة لأن العناوين
% المختلطة (عربي + إنجليزي بين أقواس) تتبعثر في وضع الـ mark مع bidi.
\fancyhead[L]{\small\color{spInk}\ifnum\value{section}>0 {\arabicfont الفصل}~\thesection\fi}
\fancyfoot[C]{\small\thepage}

% ── تنسيق الأكواد (listings) — تُستخدَم داخل بيئة english ────────────
\lstset{basicstyle=\monofont\small,backgroundcolor=\color{spGrayBg},
  frame=single,rulecolor=\color{spRule},breaklines=true,breakatwhitespace=true,
  keywordstyle=\color{spPrimary}\bfseries,commentstyle=\color{spWhy}\itshape,
  stringstyle=\color{spSec},columns=fullflexible,showstringspaces=false,
  xleftmargin=4pt,framexleftmargin=2pt,tabsize=2}
% ملاحظة: الأكواد تُكتب دائمًا داخل بيئة english يدويًّا (نمط مثبَّت مع RTL):
%   \begin{english}\begin{lstlisting}[language={[Sharp]C}] ... \end{lstlisting}\end{english}
% لا تستخدم lstnewenvironment مع english (يبتلع باقي المستند في وضع bidi).

% ── الجداول ──────────────────────────────────────────────────────────
\newcolumntype{L}[1]{>{\raggedright\arraybackslash}p{#1}}
\newcolumntype{C}[1]{>{\centering\arraybackslash}p{#1}}
\newcolumntype{Y}{>{\raggedright\arraybackslash}X}
\arrayrulecolor{spRule}
\renewcommand{\arraystretch}{1.28}
\setlength{\tabcolsep}{5pt}

% ── الصناديق الملوّنة (mdframed آمن مع RTL — لا تستخدم tcolorbox) ────
% نمط أساسي مشترك لكل الصناديق المعنونة
\global\mdfdefinestyle{spbase}{%
  linewidth=0.8pt,roundcorner=3pt,
  innerleftmargin=9pt,innerrightmargin=9pt,
  innertopmargin=7pt,innerbottommargin=7pt,
  skipabove=7pt,skipbelow=7pt,
  frametitlefont=\bfseries\color{white},
  frametitlerule=false,
  frametitleaboveskip=4pt,frametitlebelowskip=4pt}

% مولّد بيئة صندوق: #1 مفتاح لوني نصّي، #2 خلفية، #3 خلفية العنوان، #4 عنوان افتراضي
% كل بيئة تقبل عنوانًا اختياريًا لتخصيص السطر (مثل «الفكرة: RAG»)
\newenvironment{idea}[1][الفكرة]%
  {\begin{mdframed}[style=spbase,linecolor=spHow!75!black,backgroundcolor=spLightBlue,
     frametitlebackgroundcolor=spHow,frametitle={\strut #1}]}%
  {\end{mdframed}}
\newenvironment{why}[1][ليه بالطريقة دي؟]%
  {\begin{mdframed}[style=spbase,linecolor=spWhy!70!black,backgroundcolor=spLightGreen,
     frametitlebackgroundcolor=spWhy!80!black,frametitle={\strut #1}]}%
  {\end{mdframed}}
\newenvironment{keyp}[1][نقطة مهمة]%
  {\begin{mdframed}[style=spbase,linecolor=spQA!70!black,backgroundcolor=spLightPurple,
     frametitlebackgroundcolor=spQA!80!black,frametitle={\strut #1}]}%
  {\end{mdframed}}
\newenvironment{secnote}[1][زاوية الأمان]%
  {\begin{mdframed}[style=spbase,linecolor=spSec!75!black,backgroundcolor=spLightOrange,
     frametitlebackgroundcolor=spSec!80!black,frametitle={\strut #1}]}%
  {\end{mdframed}}
\newenvironment{qa}[1][سؤال متوقع في المناقشة]%
  {\begin{mdframed}[style=spbase,linecolor=spQA!70!black,backgroundcolor=spLightPurple,
     frametitlebackgroundcolor=spQA!70!black,frametitle={\strut #1}]}%
  {\end{mdframed}}
\newenvironment{note}[1][ملاحظة]%
  {\begin{mdframed}[style=spbase,linecolor=spAccent!70!black,backgroundcolor=spLightTeal,
     frametitlebackgroundcolor=spAccent!80!black,frametitle={\strut #1}]}%
  {\end{mdframed}}
% صندوق «ملخّص الأفكار بالعربي» — الأبرز؛ حدّ أعرض
\newenvironment{summary}[1][الخلاصة في سطور]%
  {\begin{mdframed}[style=spbase,linewidth=1.1pt,linecolor=spPrimary,
     backgroundcolor=spGrayBg,frametitlebackgroundcolor=spPrimary,frametitle={\strut #1}]}%
  {\end{mdframed}}
% صندوق «ركّز هنا لإبهار اللجنة» — تمييز ذهبيّ بعنوان داكن لتباين واضح
\newenvironment{spotlight}[1][ركّز هنا لإبهار اللجنة]%
  {\begin{mdframed}[style=spbase,linewidth=1.1pt,linecolor=spYellow!85!black,
     backgroundcolor=spLightYellow,frametitlebackgroundcolor=spYellow,
     frametitlefont=\bfseries\color{spInk},frametitle={\strut #1}]}%
  {\end{mdframed}}

% صندوق العنوان (الغلاف) لكل مستند
\newcommand{\coverbox}[2]{%
  \begin{mdframed}[style=spbase,linewidth=0pt,roundcorner=4pt,
     backgroundcolor=spPrimary,innertopmargin=12pt,innerbottommargin=12pt]%
  {\color{white}\Huge\bfseries ScholarPath}\\[3pt]
  {\color{white}\LARGE #1}\\[6pt]
  {\color{white}\normalsize #2}%
  \end{mdframed}}
